\documentclass{article}
\usepackage{amsmath}
\usepackage{graphicx}

\title{Physics-Informed Neural Networks for Accelerated CFD: A Benchmark Study}
\author{Samba Ba \\ PhysicsAI-Lab}
\date{\today}

\begin{document}
\maketitle

\begin{abstract}
Computational Fluid Dynamics (CFD) simulations are essential but computationally expensive. This paper explores the use of Physics-Informed Neural Networks (PINNs) as surrogates for steady and unsteady flow problems. We compare PINN predictions against high-fidelity OpenFOAM simulations for the Burgers equation and flow around a cylinder. Our results show that PINNs can achieve reasonable accuracy with up to 100x speedup, making them suitable for real-time applications and design optimization.
\end{abstract}

\section{Introduction}
CFD solves Navier-Stokes equations numerically, requiring hours to days for complex geometries. Recent advances in scientific machine learning propose using neural networks that incorporate physical laws as soft constraints. This work evaluates PINNs for surrogate modeling.

\section{Methodology}
We generate a dataset using OpenFOAM for 2D laminar flow. A PINN with 4 hidden layers and tanh activations is trained to predict velocity and pressure fields given coordinates and time. The loss function combines data mismatch and PDE residuals (Burgers' equation).

\section{Results}
Figure~\ref{fig:comparison} shows the predicted velocity field versus ground truth. Mean relative error is below 5\% on the test set, while inference time is under 0.01 seconds per point.

\begin{figure}[h]
\centering
%\includegraphics[width=0.8\textwidth]{comparison.png}
\caption{Comparison of CFD (left) and PINN prediction (right) for flow around a cylinder.}
\label{fig:comparison}
\end{figure}

\section{Conclusion}
PINNs offer a promising path toward real-time simulation. Future work will extend to turbulent flows and 3D geometries.

\bibliographystyle{plain}
\bibliography{references}
\end{document}
